\section{Trabalhos Relacionados}
\label{sec:trabalhos_relacionados}

Esta seção apresenta estudos sobre a qualidade das descrições MCP e a identificação de \textit{smells}. A literatura mapeia desafios de manutenção e vulnerabilidades que comprometem o desempenho de agentes autônomos. A análise desses trabalhos fundamenta a lacuna motivadora desta pesquisa. Busca-se superar a avaliação isolada de documentações por meio da integração do código, garantindo a fidelidade entre o texto e a implementação real.

Hasan et al. (2026) \nocite{2026arXiv260214878M} investigaram a qualidade das documentações de 856 ferramentas no ecossistema MCP por meio de uma análise empírica fundamentada em LLMs. O estudo aplicou uma rubrica de seis dimensões focada na clareza e completude do conteúdo textual. Os resultados revelaram que 97,1\% das descrições possuem ao menos um \textit{smell}, o que reduz a precisão dos agentes na seleção de parâmetros. \textcolor{red}{O presente estudo parte dessa investigação e avança em uma direção não explorada pelos autores, sendo ela a incorporação do
código ao processo de avaliação. Enquanto a abordagem anterior analisa as descrições de forma isolada, este trabalho utiliza o contexto técnico, extraído via grafos de chamadas, para confrontar a intenção declarada com a implementação real, verificando em que medida o acesso à
lógica interna do servidor altera os escores obtidos nas dimensões da rubrica proposta pelos autores}.

Em seu estudo sobre inconsistências entre descrições e implementações de ferramentas MCP, Shi et al. (2026) \nocite{shi2026description} investigam a prevalência e os impactos de inconsistências semânticas entre descrições de ferramentas e seus respectivos códigos. O método consistiu em propor uma taxonomia de inconsistências descrição-código e um detector baseado em LLM para analisar pares de descrição e implementação em larga escala em servidores MCP reais. Os resultados revelaram que 9,93\% das ferramentas analisadas apresentam inconsistências e que 35\% dos servidores possuem ao menos um caso desse problema, evidenciando riscos relacionados à seleção incorreta de ferramentas, efeitos colaterais não declarados e falhas de segurança. \textcolor{red}{Esse estudo sustenta a premissa desta pesquisa ao demonstrar que informações relevantes para compreender o comportamento de uma ferramenta frequentemente não estão presentes em sua descrição, o que fornece evidências de que a análise do código pode enriquecer a avaliação da qualidade de descrições MCP realizada por LLMs.}

O panorama de adoção e maturidade do ecossistema MCP é mapeado empiricamente por Guo et al. (2025)\nocite{guo2025measurement}. O trabalho implementa um rastreador sistemático que analisa mais de 8.400 projetos em grandes repositórios para caracterizar o estado real da adoção do protocolo. Entre os achados, os autores constatam que mais da metade dos projetos listados são inválidos, abandonados ou de baixo valor agregado, o que evidencia a heterogeneidade de qualidade do ecossistema. Esse estudo fornece uma visão macro do MCP, mas concentra-se na caracterização dos repositórios e não examina a consistência entre as descrições das ferramentas e sua implementação no código. Este trabalho situa-se nesse ponto não coberto, ao investigar a qualidade das descrições no nível da ferramenta e o papel do código em sua avaliação. O dado de baixa maturidade reportado pelos autores também fundamenta o critério de popularidade adotado na seleção da amostra, descrito na Seção~\ref{sec:materiais}.

Com foco no desempenho dos agentes, Luo et al. (2025)\nocite{2025arXiv250814704L} propuseram o MCP-Universe, um \textit{benchmark} para avaliar modelos de linguagem em servidores MCP reais. O estudo emprega avaliadores baseados em execução, que validam o resultado prático das chamadas de ferramentas em vez de analisar apenas a semelhança textual das respostas. Ao avaliar mais de 20 modelos em 231 tarefas distribuídas por 11 servidores, os resultados revelaram que mesmo modelos de fronteira atingem taxas de sucesso abaixo de 45\%, com falhas frequentemente associadas ao uso incorreto de interfaces de ferramentas pouco familiares. O presente trabalho investiga uma possível causa raiz dessas falhas, verificando se a divergência semântica entre a descrição textual de uma ferramenta e seu comportamento real é o fator que induz os agentes ao erro nesses cenários.

Em seu estudo sobre avaliação automática de modelos de linguagem, Zheng et al. (2023) \nocite{NEURIPS2023_91f18a12} investigam a viabilidade da abordagem \textit{LLM-as-a-Judge}, na qual modelos de linguagem são utilizados para avaliar a qualidade das respostas produzidas por outros modelos. O método consistiu na comparação entre julgamentos realizados por GPT-4 e avaliações humanas em \textit{benchmarks} de conversação. Os resultados mostraram que o GPT-4 alcançou concordância superior a 80\% com avaliadores humanos, desempenho comparável ao observado entre os próprios especialistas. Este trabalho fundamenta metodologicamente a pesquisa proposta ao fornecer evidências da confiabilidade de avaliadores baseados em LLM, empregados neste estudo para analisar a qualidade de descrições MCP e o impacto da inclusão do código nesse processo de avaliação.

Em síntese, e demonstrado que a qualidade das descrições MCP é um fator crítico, pois discrepâncias entre o texto e a implementação comprometem a eficácia dos agentes. Embora estudos validem o uso de LLMs como avaliadores e a importância do código para a compreensão sistêmica, as abordagens atuais analisam as documentações de forma isolada. Diante dessa lacuna, este trabalho investiga como a incorporação do contexto técnico influencia a detecção de \textit{smells} em descrições MCP.
